\documentclass[10pt]{beamer}
\usetheme{metropolis}
\usepackage{graphicx}
\usepackage{listings}
\usepackage{booktabs}
\usepackage{xcolor}
\definecolor{codebg}{HTML}{F5F5F5}
\definecolor{codegreen}{HTML}{2E7D32}
\definecolor{codegray}{HTML}{757575}
\definecolor{codeblue}{HTML}{1565C0}
\lstdefinestyle{rstyle}{language=R,backgroundcolor=\color{codebg},basicstyle=\ttfamily\scriptsize,
  keywordstyle=\color{codeblue}\bfseries,stringstyle=\color{codegreen},
  commentstyle=\color{codegray}\itshape,breaklines=true,frame=single,
  rulecolor=\color{codegray!30},showstringspaces=false,tabsize=2,
  morekeywords={library,skim,skimr,create_report,DataExplorer,dfSummary,
    summarytools,render,quarto,rmarkdown,params,knit}}
\lstdefinestyle{pythonstyle}{language=Python,backgroundcolor=\color{codebg},basicstyle=\ttfamily\scriptsize,
  keywordstyle=\color{codeblue}\bfseries,stringstyle=\color{codegreen},
  commentstyle=\color{codegray}\itshape,breaklines=true,frame=single,
  rulecolor=\color{codegray!30},showstringspaces=false,tabsize=4,
  morekeywords={import,from,as,pd,ProfileReport,ydata_profiling,
    sweetviz,analyze,compare,dtale,show,to_file,to_notebook_iframe}}
\lstdefinestyle{yaml}{backgroundcolor=\color{codebg},basicstyle=\ttfamily\scriptsize,
  keywordstyle=\color{codeblue}\bfseries,commentstyle=\color{codegray}\itshape,
  breaklines=true,frame=single,rulecolor=\color{codegray!30},showstringspaces=false}
\title{EDA Automation \& Reproducible Reports}
\subtitle{Workshop 4 --- Module 9}
\author{Prof.Asoc.\ Endri Raco}
\institute{Data Visualization for Data Scientists \\ UNYT Tirana}
\date{2025/26}
\begin{document}

\begin{frame}
  \titlepage
\end{frame}

\begin{frame}{Module Roadmap}
  \begin{enumerate}
    \item Manual vs automated EDA: complementary roles
    \item Auto-EDA packages: skimr, DataExplorer, ydata-profiling
    \item The hybrid workflow: auto baseline $\rightarrow$ manual deep-dive
    \item Reproducible reports: Quarto, RMarkdown, Jupyter
    \item Parameterised reports: one template, many outputs
  \end{enumerate}
  \begin{exampleblock}{Learning Outcome}
    You will generate a comprehensive automated EDA report in under 60
    seconds, identify which findings require manual follow-up, and
    produce a reproducible report that combines code, prose, and output
    in a single document.
  \end{exampleblock}
\end{frame}

\section{Manual vs Automated}

\begin{frame}{Two Approaches, One Goal}
  \begin{center}
    \includegraphics[width=0.92\textwidth]{figures_w04m09/manual_vs_auto}
  \end{center}
\end{frame}

\begin{frame}{When to Use Which}
  \centering
  \begin{tabular}{lll}
    \toprule
    \textbf{Situation} & \textbf{Approach} & \textbf{Why} \\
    \midrule
    First look at unknown data & Automated & Fast, comprehensive baseline \\
    Client presentation & Manual & Full narrative control \\
    20 datasets to profile & Automated & Batch processing in a loop \\
    Deep-dive on one finding & Manual & Custom chart for the story \\
    Coursework / publication & Both & Auto for audit, manual for report \\
    \bottomrule
  \end{tabular}
  \vspace{6pt}
  \begin{alertblock}{Pro-Tip}
    Auto-EDA is a \textbf{screening tool}, not a replacement for thinking.
    It tells you \emph{what to look at}; you decide \emph{what it means}.
    The hybrid workflow uses auto-EDA for the first 30 seconds, then
    spends 2 hours on the findings that matter.
  \end{alertblock}
\end{frame}

\section{Auto-EDA Packages}

\begin{frame}{The Package Landscape}
  \begin{center}
    \includegraphics[width=0.95\textwidth]{figures_w04m09/landscape}
  \end{center}
\end{frame}

\begin{frame}[fragile]{R: skimr + DataExplorer}
\begin{lstlisting}[style=rstyle]
# ── skimr: compact console summary ──────────────
library(skimr)
skim(apps)                  # full summary
skim(apps) |> summary()     # just counts

# Grouped skim (like group_by + summarise but for EDA)
apps |> group_by(Type) |> skim(Rating, Reviews)

# ── DataExplorer: one-click HTML report ─────────
library(DataExplorer)
create_report(apps)         # full HTML report
# Or individual plots:
plot_missing(apps)          # missingness bar
plot_histogram(apps)        # all numeric histograms
plot_correlation(apps)      # correlation heatmap
plot_bar(apps)              # all categorical bars

# ── summarytools: freq tables + embedded graphs ──
library(summarytools)
dfSummary(apps) |> view()   # opens in browser
\end{lstlisting}
\end{frame}

\begin{frame}[fragile]{Python: ydata-profiling + sweetviz}
\begin{lstlisting}[style=pythonstyle]
# ── ydata-profiling (formerly pandas-profiling) ──
# pip install ydata-profiling
from ydata_profiling import ProfileReport

report = ProfileReport(
    apps,
    title="Google Play Store EDA",
    explorative=True,           # interactions tab
    correlations={
        "pearson": True,
        "spearman": True,
        "kendall": False})

report.to_file("eda_report.html")       # standalone HTML
report.to_notebook_iframe()              # inline in Jupyter

# ── sweetviz: compare two datasets ──────────────
# pip install sweetviz
import sweetviz as sv

# Compare Free vs Paid
report = sv.compare(
    [apps.query("Type=='Free'"), "Free"],
    [apps.query("Type=='Paid'"), "Paid"])
report.show_html("free_vs_paid.html")
\end{lstlisting}
\end{frame}

\section{The Hybrid Workflow}

\begin{frame}{Auto Baseline $\rightarrow$ Manual Deep-Dive}
  \begin{center}
    \includegraphics[width=0.95\textwidth]{figures_w04m09/workflow}
  \end{center}
\end{frame}

\begin{frame}{What to Look for in the Auto Report}
  \centering
  \begin{tabular}{llp{5cm}}
    \toprule
    \textbf{Section} & \textbf{Scan For} & \textbf{Follow-Up Action} \\
    \midrule
    Overview & Row/col counts, types & Verify expected dimensions \\
    Alerts & High correlation, skew, zeros, constants & Investigate each alert \\
    Missing & \% per column, patterns & Shadow analysis if $>$5\% \\
    Distributions & Skew, bimodality, spikes & Manual histogram with annotation \\
    Correlations & $|r| > 0.7$ & Scatter plot for top pairs \\
    Interactions & Unexpected nonlinearities & Faceted scatter + smoother \\
    \bottomrule
  \end{tabular}
  \vspace{6pt}
  {\small The auto report generates $\sim$50 plots in 30 seconds. Your
  job: scan for the 3--5 that require a custom, publication-quality
  deep-dive.}
\end{frame}

\section{Reproducible Reports}

\begin{frame}{Code + Prose + Output = One Document}
  \begin{center}
    \includegraphics[width=0.95\textwidth]{figures_w04m09/repro_tools}
  \end{center}
\end{frame}

\begin{frame}[fragile]{Quarto: The Modern Standard}
\begin{lstlisting}[style=yaml]
---
title: "EDA Report: Google Play Store"
author: "Endri Raco"
format:
  html:
    toc: true
    code-fold: true
    theme: cosmo
  pdf:
    documentclass: article
execute:
  echo: true
  warning: false
---

## Data Cleaning
```{r}
library(tidyverse)
apps <- read_csv("googleplaystore.csv") |>
  distinct(App, .keep_all = TRUE) |>
  mutate(Reviews = as.numeric(Reviews))
```

## Distribution
```{r}
#| fig-cap: "Rating distribution (left-skewed)"
ggplot(apps, aes(x = Rating)) + geom_histogram(bins = 30)
```

## Key Finding
Rating is **left-skewed** (mean = `r round(mean(apps$Rating, na.rm=T), 2)`).
\end{lstlisting}
\end{frame}

\begin{frame}[fragile]{Jupyter Notebook: Python Standard}
\begin{lstlisting}[style=pythonstyle]
# Cell 1: Markdown
# # EDA Report: Google Play Store
# **Author**: Endri Raco | **Date**: 2025

# Cell 2: Code
import pandas as pd
apps = pd.read_csv("googleplaystore.csv")
apps.describe()

# Cell 3: Visualization
apps["Rating"].hist(bins=30, color="#1565C0")
plt.title("Rating Distribution (left-skewed)")
plt.show()

# Cell 4: Markdown
# ## Key Finding
# Rating is left-skewed (mean = 4.17, median = 4.30).

# Export:
# jupyter nbconvert --to pdf notebook.ipynb
# jupyter nbconvert --to html notebook.ipynb
# quarto render notebook.ipynb --to pdf
\end{lstlisting}
\end{frame}

\section{Parameterised Reports}

\begin{frame}{One Template, Many Outputs}
  \begin{center}
    \includegraphics[width=0.88\textwidth]{figures_w04m09/parameterised}
  \end{center}
\end{frame}

\begin{frame}[fragile]{Parameterised Quarto / RMarkdown}
\begin{lstlisting}[style=yaml]
---
title: "EDA Report"
params:
  dataset: "googleplaystore.csv"
  target: "Rating"
  group: "Type"
---
```{r}
df <- read_csv(params$dataset) |> drop_na(!!sym(params$target))
ggplot(df, aes(x = .data[[params$target]])) +
  geom_histogram(bins = 30) +
  facet_wrap(reformulate(params$group))
```
\end{lstlisting}
  \vspace{4pt}
\begin{lstlisting}[style=rstyle]
# Render for each dataset:
quarto::quarto_render("template.qmd",
  execute_params = list(dataset = "netflix.csv",
    target = "duration", group = "type"))

# Batch render all coursework variants:
for (v in 0:9) {
  quarto::quarto_render("template.qmd",
    output_file = paste0("report_v", v, ".html"),
    execute_params = list(variant = v))
}
\end{lstlisting}
\end{frame}

\begin{frame}{Why Reproducibility Matters}
  \centering
  \begin{tabular}{lp{5.5cm}}
    \toprule
    \textbf{Benefit} & \textbf{How} \\
    \midrule
    Auditability & Reviewer re-runs your code, gets identical output \\
    Version control & \texttt{git diff} shows exactly what changed \\
    Error correction & Fix one line, re-render; no manual copy-paste \\
    Scalability & Same report for 10 datasets via parameterisation \\
    Collaboration & Co-authors edit code, not screenshots \\
    Publication & Journals increasingly require reproducible code \\
    \bottomrule
  \end{tabular}
  \vspace{6pt}
  \begin{alertblock}{Pro-Tip}
    The coursework requires R \emph{and} Python scripts plus a report.
    A Quarto document with both \texttt{\{r\}} and \texttt{\{python\}}
    code chunks can produce the entire submission as \textbf{one
    reproducible file}. The \texttt{reticulate} package bridges R and
    Python in the same document.
  \end{alertblock}
\end{frame}

\section{Synthesis}

\begin{frame}{Key Takeaways}
  \begin{enumerate}
    \item \textbf{Auto-EDA is a screening tool}: 30-second baseline that
          tells you what to investigate manually.
    \item In R: \texttt{skimr::skim()} for console, \texttt{DataExplorer::
          create\_report()} for HTML. In Python: \texttt{ydata-profiling}
          for comprehensive reports, \texttt{sweetviz} for comparisons.
    \item The \textbf{hybrid workflow}: auto report $\rightarrow$ scan
          for anomalies $\rightarrow$ manual deep-dive $\rightarrow$
          compose dashboard $\rightarrow$ write narrative.
    \item \textbf{Quarto} (.qmd) is the modern standard for reproducible
          reports: code + prose + output in one file, renders to HTML,
          PDF, slides, dashboards.
    \item \textbf{Parameterised reports} generate many outputs from one
          template --- ideal for batch EDA and coursework grading.
    \item \textbf{Reproducibility} = auditability + version control +
          error correction + scalability.
  \end{enumerate}
\end{frame}

\begin{frame}{Essential Readings}
  \begin{itemize}
    \item Quarto documentation: \texttt{https://quarto.org}
    \item Wickham, H. \& Grolemund, G. (2023). \emph{R for Data Science},
          2nd ed., Chapters 28--29 (Quarto).
    \item ydata-profiling: \texttt{https://docs.profiling.ydata.ai}
    \item Peng, R. (2011). ``Reproducible Research in Computational
          Science.'' \emph{Science}, 334(6060).
  \end{itemize}
  \vspace{6pt}
  \textbf{Next Module}: Lab --- Full EDA on Netflix (W04-M10)
\end{frame}

\end{document}
